\subsection{Adjoint of a homomorphism.}

In this section, we denote by A a ring (resp. a ring equipped with an antiautomorphism J), by E and E' two left A-modules, by F and F' two right A-modules (resp. left A-modules), and by Φ and Φ' two bilinear forms (resp. sesquilinear forms for J) on E × F and E' × F' respectively. We assume that Φ is nondegenerate, in other words (no. 1) that the linear maps \(d_{\Phi}\) and \(s_{\Phi}\) associated with Φ are injective.

Given a homomorphism \(u\) from E into \(\mathbf{E}'\) , consider the set \(\mathbf{F}_{1}'\) of elements \(y'\) of \(\mathbf{F}'\) such that there exists \(y \in \mathbf{F}\) for which one has \(d_{\Phi'}(y') \circ u = d_{\Phi}(y)\) , that is to say \(\Phi'(u(x), y') = \Phi(x, y)\) for all \(x \in \mathbf{E}\) . It is clear that \(\mathbf{F}_{1}'\) is a submodule of \(\mathbf{F}'\) . Since \(d_{\Phi}\) is injective, there exists, for every \(y' \in \mathbf{F}_{1}'\) , one and only one element \(y\) of F such that \(\Phi'(u(x), y') = \Phi(x, y)\) . The map \(y' \to y\) of \(\mathbf{F}_{1}'\) into F thus defined is A-linear; denoting it by \(u^{*}\) , one has, for every \(x \in \mathbf{E}\) and every \(y' \in \mathbf{F}_{1}'\)

\[
\Phi^ {\prime} (u (x), y ^ {\prime}) = \Phi (x, u ^ {*} (y ^ {\prime})).\tag{30}
\]

DEFINITION 10. --- The hypotheses and notations being as before, one says that the homomorphism \(u^*\) of \(F_1'\) into F satisfying (30) is the left adjoint of \(u\) , and that \(F_1'\) is the defining submodule of \(u^*\) .

One defines similarly the right adjoint of a homomorphism \(\nu\) from F into F' by the formula

\[
\Phi^ {\prime} (x ^ {\prime}, \nu (y)) = \Phi (\nu^ {*} (x ^ {\prime}), y) \quad (x ^ {\prime} \in \mathrm{E} _ {1} ^ {\prime}, y \in \mathrm{F}),\tag{31}
\]

where \(E_{1}^{\prime}\) denotes the submodule of \(E^{\prime}\) defined analogously to \(F_{1}^{\prime}\) .

Remark. --- If the left adjoint \(u^{*}\) of \(u: E \to E'\) is everywhere defined, and if \(s_{\Phi'}\) and \(d_{\Phi'}\) are injective, formula (30) shows that \(u\) is the right adjoint of \(u^{*}\) .

From (30), we deduce that, if \(u_1\) and \(u_2\) are two homomorphisms from E into E′ admitting everywhere-defined adjoints, and if c is an element of the center of A, then we have

\[
\left\{ \begin{array}{l} (u _ {1} + u _ {2}) ^ {*} = u _ {1} ^ {*} + u _ {2} ^ {*}; 1 ^ {*} = 1; \\ (c u _ {1}) ^ {*} = c. u _ {1} ^ {*} \quad \text { when } \Phi \text { and } \Phi^ {\prime} \text { are bilinear }; \\ (c u _ {1}) ^ {*} = c ^ {J ^ {\prime}}. u _ {1} ^ {*} \quad \text { when } \Phi \text { and } \Phi^ {\prime} \text { are sesquilinear }. \end{array} \right.\tag{32}
\]

Additionally, if \(E^{\prime\prime}\) is a third left A-module, \(F^{\prime\prime}\) a third right (resp. left) A-module, \(\Phi^{\prime\prime}\) a bilinear form (resp. sesquilinear form with respect to J) on \(E^{\prime\prime} \times F^{\prime\prime}\) , and if \(u^{\prime}\) is a homomorphism of \(E^{\prime}\) into \(E^{\prime\prime}\) admitting a (left) adjoint defined everywhere, we have

\[
(u ^ {\prime} \circ u) ^ {*} = u ^ {*} \circ u ^ {\prime *}.\tag{33}
\]

In particular, if \(u\) is an isomorphism from E onto E', and if the adjoints \(u^{*}\) and \((u^{-1})^{*}\) are everywhere defined, \(u^{*}\) is an isomorphism from F' onto F, and we have \((u^{*})^{-1} = (u^{-1})^{*}\) . Analogous properties hold for right adjoints.

PROPOSITION 7. --- With the same notation as above, assume that \(d_{\Phi}\) is bijective. Then every homomorphism \(u\) from E to E' admits an everywhere-defined left adjoint, and we have \(u^{*} = (d_{\Phi})^{-1} \circ {}^{t}u \circ d_{\Phi'}\) .

Indeed, as \(d_{\Phi}\) is bijective, we have, with the notations from the beginning of the paragraph, \(\mathbf{F}_{1}^{\prime}=\mathbf{F}^{\prime}\) , and \(u^{*}\) is therefore everywhere defined. On the other hand, (30) is equivalent to

\[
\left\langle u (x), d _ {\Phi^ {\prime}} \left(y ^ {\prime}\right) \right\rangle = \left\langle x, \left(d _ {\Phi} \circ u ^ {*}\right) \left(y ^ {\prime}\right) \right\rangle \quad \left(x \in \mathrm{E}, y ^ {\prime} \in \mathrm{F} ^ {\prime}\right);
\]

or \(\langle d_{\Phi'}(y'), u(x) \rangle = \langle {}^t u(d_{\Phi'}(y')), x \rangle\) ; one therefore has \({}^t u(d_{\Phi'}(y')) = d_{\Phi}(u^*(y'))\) for all \(y' \in F'\) , whence \({}^t u \circ d_{\Phi'} = d_{\Phi} \circ u^*\) , and consequently the announced expression of \(u^*\) . QED.

Remark. --- When \(s_{\Phi}\) is bijective, every homomorphism \(\nu\) from F to F' admits a right adjoint that is everywhere defined, and we have

\[
v ^ {*} = (s _ {\Phi}) ^ {- 1} \circ {} ^ {t} v \circ s _ {\Phi^ {\prime}}.\tag{34}
\]

PROPOSITION 8. --- With the same notation as above, assume that \(s_{\Phi}\) and \(d_{\Phi}\) are bijective. Let \(u\) and \(v\) be isomorphisms from E onto \(\mathbf{E}'\) and from F onto \(\mathbf{F}'\) respectively. Then, for \(\Phi\) to be the inverse image of \(\Phi'\) with respect to \(u\) and \(\nu\) (that is, so that one has \(\Phi(x,y)=\Phi'(u(x),\nu(y))\) for all \(x\in\mathbf{E},y\in\mathbf{F}\) ), it is necessary and sufficient that we have \(u^{-1}=\nu^{*}\) and \(\nu^{-1}=u^{*}\) .

Indeed \(\Phi'(u(x), v(y)) = \Phi(x, y)\) is also written \(\Phi(x, u^*(v(y))) = \Phi(x, y)\) . If this holds for all \(x \in E\) and \(y \in F\) , one has \(u^* \circ v = 1\) since \(\Phi\) is nondegenerate. We therefore also have \(v^* \circ u = 1\) by (33). The converse is immediate.

COROLLARY. --- Let A be a ring equipped with an antiautomorphism J, and let E be a left A-module, \(\Phi\) a sesquilinear form for J on E × E whose associated maps are bijective, and u an automorphism of the A-module E. For u to leave \(\Phi\) invariant (that is, so that we have \(\Phi(u(x), u(y)) = \Phi(x, y)\) for all x, y in E), it is necessary and sufficient that the two adjoints of u be equal and that we have \(u^{*} = u^{-1}\) .

This follows immediately from prop. 8.

Remark. --- Under the hypotheses of the corollary to Proposition 8, suppose further that A is a field and that E is finite-dimensional over A. Let \(w\) be an endomorphism of E, \(w_{1}\) and \(w_{2}\) its right and left adjoints. Each of the conditions \(ww_{1}=1\) , \(ww_{2}=1\) , \(w_{1}w=1\) , \(w_{2}w=1\) implies that \(w\) is an automorphism of E leaving \(\Phi\) invariant, and that \(w_{1}=w_{2}\) .

